\bookchapter{Appendix A: New Phone Unboxing “Step-by-Step” Guide}

\booksection{Unboxing Inspection “Three-Step Process”:}

\begin{itemize}
  \item \textbf{: Make a test call to check the earpiece and microphone, take a photo to check the camera.}2. Activation \& Initial Setup "Key Reminders":
  \item \textbf{Network}: Connect to a stable Wi-Fi first.
  \item \textbf{Account}: Log in to your Apple ID, Huawei account, Xiaomi account, etc.
\end{itemize}

\booksection{Activation \& Initial Setup “Key Reminders”:}

\begin{itemize}
  \item \textbf{: Turn on "Find My Phone" and set a reliable lock screen password.}eSIM Setup
  \item \textbf{: If your phone supports eSIM (e.g., iPhone, some Android flagships), you can add an eSIM plan in Settings without inserting a physical SIM card. eSIM is especially convenient when traveling abroad—you can directly purchase a local carrier's eSIM plan.}Data Migration
  \item \textbf{: When migrating data from an old phone, ensure both phones have sufficient battery. iPhone uses "Quick Start", Android uses the official device transfer app.}3. Installing Apps "Safety First":
  \item \textbf{Use Official Stores First}: Apple uses App Store, Android uses the phone's built-in app store.
  \item \textbf{Don't Grant Permissions Carelessly}: When permission requests pop up during installation, think carefully: "Does this app really need this?"
\end{itemize}

\booksection{Installing Apps “Safety First”:}

\begin{itemize}
  \item \textbf{Basic Practice}: Clean up phone storage—the goal is to find and delete those old videos you "forgot what they were", duplicate screenshots, and apps you haven't used in a long time.
  \item \textbf{Software Practice}: Set up WeChat storage cleanup. The path is usually [Me] - [Settings] - [General] - [Storage].
\end{itemize}
